\documentclass[11pt]{article}
\usepackage[margin=1.1in]{geometry}
\usepackage[T1]{fontenc}
\usepackage{lmodern}
\usepackage{microtype}
\usepackage{amsmath,amssymb,amsthm}
\usepackage{mathtools}
\usepackage[most]{tcolorbox}
\usepackage{listings}
\usepackage{enumitem}
\usepackage{booktabs}
\usepackage{array}
\usepackage[colorlinks=true,linkcolor=black,urlcolor=blue!60!black]{hyperref}

% ---------- theorem machinery (numbered within each lesson) ----------
\newtheorem{theorem}{Theorem}[section]
\newtheorem{proposition}[theorem]{Proposition}
\newtheorem{lemma}[theorem]{Lemma}
\newtheorem{corollary}[theorem]{Corollary}
\theoremstyle{definition}
\newtheorem{definition}[theorem]{Definition}
\newtheorem{example}[theorem]{Example}
\newtheorem{remark}[theorem]{Remark}
\newtheorem{innerexercise}{Exercise}[section]
\newenvironment{exercise}[1][]{\begin{innerexercise}[#1]}{\end{innerexercise}}

% ---------- notation ----------
\newcommand{\R}{\mathbb{R}}
\newcommand{\C}{\mathbb{C}}
\newcommand{\Z}{\mathbb{Z}}
\newcommand{\ip}[2]{\langle #1,\, #2\rangle}
\newcommand{\norm}[1]{\lVert #1\rVert}
\newcommand{\abs}[1]{\lvert #1\rvert}
\newcommand{\one}{\mathbf{1}}
\newcommand{\diag}{\operatorname{diag}}
\newcommand{\median}{\operatorname{median}}
\newcommand{\round}{\operatorname{round}}
\newcommand{\cconv}{\circledast}
\newcommand{\code}[1]{\texttt{#1}}

% ---------- "used in" tag: where each lesson lives in the code ----------
\newcommand{\usedin}[1]{\par\noindent\colorbox{black!8}{\parbox{\dimexpr\linewidth-2\fboxsep\relax}{\small\textsc{Used in:}~#1}}\par\medskip}

% ---------- boxes ----------
\newtcolorbox{connection}{colback=black!4,colframe=black!55,boxrule=0.5pt,
  arc=1.5pt,left=7pt,right=7pt,top=5pt,bottom=5pt,breakable,
  before skip=10pt,after skip=10pt}

% ---------- code ----------
\lstset{
  language=Python,
  basicstyle=\ttfamily\small,
  commentstyle=\itshape\color{black!60},
  keywordstyle=\bfseries,
  showstringspaces=false,
  columns=fullflexible,
  keepspaces=true,
  frame=single,
  framerule=0.3pt,
  rulecolor=\color{black!40},
  xleftmargin=6pt,xrightmargin=6pt,
  aboveskip=8pt,belowskip=8pt,
  breaklines=true
}

\setlist{itemsep=2pt,topsep=4pt}
\setcounter{tocdepth}{1}

\title{\bfseries The Mathematics of Humming a Song\\[4pt]
\large A short primer for \texttt{ocarina-hum-matcher}}
\author{Diivanand Ramalingam}
\date{November 2026}

\begin{document}
\maketitle

\section*{How to use this booklet}

The \code{ocarina} tool listens to a few seconds of humming and names the
\emph{Ocarina of Time} song it heard. Everything it does rests on a short chain of
ideas: sampling, complex exponentials, the discrete Fourier transform as a change of
basis, how NumPy lays out its output, shifts and circulant matrices, leakage and
windows, the FFT, the harmonic product spectrum, a circular mean, and a median. This
booklet covers exactly that chain, in ten short lessons, and nothing else.

You need a first course in linear algebra (vectors, dot products, matrices, bases,
orthogonality) and comfort with complex numbers at the level of $i^2=-1$. No signal
processing is assumed: every term is defined before it is used.

Each lesson has the same three layers:
\begin{itemize}
  \item \textbf{Theory}: definitions, then theorems with complete proofs. Knowing
  \emph{why} a step works is what lets you predict when it will fail.
  \item \textbf{Practice by hand}: small examples ($N=4$ or $N=8$, six-note melodies)
  computed with pencil and paper.
  \item \textbf{Code connection}: a box pointing at the lines of
  \nolinkurl{src/ocarina_hum_matcher/recognizer.py} that implement the lesson.
\end{itemize}

Each lesson ends with two or three exercises, tagged \textbf{(Hand)} or
\textbf{(Proof)}: usually one by hand and one or two proofs. Do them all. The starred
proof in each lesson is the one that tests whether you understand \emph{why} the method
works, not just how. Appendix~\ref{app:hints} has answers and proof hints. Each lesson
takes 30 to 45 minutes, so the whole primer fits in a week of evenings.

\subsection*{Map of the booklet}

\begin{center}
\small
\begin{tabular}{@{}l l l l@{}}
\toprule
Lesson & Idea & Where in the code & Slides\\
\midrule
\ref{sec:sound} & Sampling, frequency, pitch, notes & \code{synth.py}, \code{songs.py} & 2--4\\
\ref{sec:complex} & Complex exponentials, probes & (the idea behind every bin) & 6\\
\ref{sec:dft} & The DFT as a change of basis & \code{magnitude\_spectrum} & 7\\
\ref{sec:layout} & Periodicity, real signals, \code{rfft} & \code{magnitude\_spectrum} & 8, 10\\
\ref{sec:circ} & Shifts, circulants, convolution & why only $\abs{X}$ is kept & 9\\
\ref{sec:window} & Leakage, windows, zero-padding & \code{np.hanning}, \code{zero\_pad} & 11, 14\\
\ref{sec:fft} & The FFT & \code{scipy.fft.rfft} & 12\\
\ref{sec:pitch} & Harmonic product, peak refinement & \code{harmonic\_product}, \code{refine\_peaks} & 13--14\\
\ref{sec:notes} & MIDI numbers, circular mean & \code{estimate\_detune}, \code{segment\_notes} & 15\\
\ref{sec:match} & Key-invariant matching & \code{rank\_songs}, \code{match\_song} & 16--17\\
\bottomrule
\end{tabular}
\end{center}

\noindent ``Slides'' refers to the deck that \code{ocarina slides} opens; the two are
meant to be used together.

\subsection*{Conventions}

Vectors have $N$ entries indexed $n=0,\dots,N-1$, and indices are read modulo $N$
whenever they leave that range. Throughout,
\[
\omega = e^{-2\pi i/N},\qquad (e_k)_n = e^{2\pi i kn/N}=\omega^{-kn},\qquad
X_k=\sum_{n=0}^{N-1}x_n\,\omega^{kn}.
\]
This is NumPy's convention: \code{np.fft.fft(x)[k]} is $X_k$.

\tableofcontents
\clearpage

% =====================================================================
\section{Sound as Numbers}\label{sec:sound}
\usedin{\code{synth.py} (how test hums are built), \code{songs.py} (notes as MIDI
numbers), \code{DetectorConfig} (40 ms frames, 80--1{,}200 Hz). Slides 2--4.}

A microphone turns air pressure into a voltage, and a sound card measures that voltage
at regular instants. This lesson fixes the vocabulary: what is recorded, what
``frequency'' and ``pitch'' mean, and why sampling makes some frequencies impossible to
tell apart.

\subsection{Theory}

\begin{definition}[Signal, sample, sample rate]
A \emph{signal} is a function $x:\R\to\R$; here $x(t)$ is the air pressure at the
microphone at time $t$ seconds. \emph{Sampling} at rate $f_s$ (samples per second)
records the numbers $x_n=x(n/f_s)$ for integers $n$. A recording of $N$ samples is a
vector $(x_0,\dots,x_{N-1})\in\R^N$. Describing a signal by its values at each instant
is called the \emph{time domain}.
\end{definition}

\begin{definition}[Period, frequency]
A signal is \emph{periodic} with period $T>0$ if $x(t+T)=x(t)$ for all $t$. The
smallest such $T$, when it exists, is the \emph{fundamental period}, and its
reciprocal $f=1/T$ is the \emph{frequency}, measured in hertz (Hz, repetitions per
second).
\end{definition}

\begin{definition}[Sinusoid]
A \emph{sinusoid} is $a\cos(2\pi f t+\varphi)$ with \emph{amplitude} $a\ge0$,
\emph{frequency} $f>0$ and \emph{phase} $\varphi\in\R$. Its fundamental period is
$1/f$. The phase only slides it along the time axis.
\end{definition}

A voice is not a single sinusoid. A sustained hum is well modelled by a sum of
sinusoids at whole-number multiples of one frequency.

\begin{proposition}[Harmonic sums repeat at the fundamental]\label{prop:harmperiod}
Let $f_0>0$ and
\[
x(t)=\sum_{h=1}^{H}a_h\cos(2\pi h f_0 t+\varphi_h).
\]
Then $1/f_0$ is a period of $x$. More generally, if an integer $g\ge1$ divides every
$h$ with $a_h\neq0$, then $1/(gf_0)$ is a period of $x$.
\end{proposition}
\begin{proof}
Replace $t$ by $t+1/(gf_0)$. The $h$-th term becomes
$a_h\cos(2\pi hf_0t+2\pi h/g+\varphi_h)$. When $a_h\ne0$, $h/g$ is an integer, so the
added angle is a multiple of $2\pi$ and the term is unchanged; when $a_h=0$ the term is
zero anyway. Taking $g=1$ gives the first claim.
\end{proof}

\begin{definition}[Fundamental, harmonics, pitch, timbre]
In Proposition~\ref{prop:harmperiod}, the components at $hf_0$ are the
\emph{harmonics} and $f_0$ is the \emph{fundamental frequency}. The \emph{pitch} of a
sound is how high or low it is heard; for sounds of this form the ear hears $f_0$
(when the harmonics present have no common factor), and \textbf{in this booklet
``pitch'' means $f_0$ in Hz}. The relative amplitudes $a_h$ set the \emph{timbre}: they
are why a hum and an ocarina playing the same note sound different.
\end{definition}

Pitch is multiplicative: equal musical steps are equal \emph{ratios} of frequency.

\begin{definition}[Notes and MIDI numbers]
In twelve-tone equal temperament, the note with \emph{MIDI number} $m\in\Z$ has
frequency
\[
f(m)=440\cdot 2^{(m-69)/12}\ \text{Hz}.
\]
So $m=69$ is A4 (440 Hz), $m=60$ is middle C (C4), and $m=62$ is D4. Adding 1 to $m$
is one \emph{semitone} (a frequency ratio of $2^{1/12}\approx1.0595$); adding 12 is
one \emph{octave} (a ratio of exactly 2).
\end{definition}

Sampling throws information away: between two samples, anything could happen. The
precise loss is this.

\begin{proposition}[Aliasing]\label{prop:alias}
Let $f_s>0$. For every integer $j$ and all integers $n$,
\[
\cos\!\Big(2\pi\frac{f}{f_s}n+\varphi\Big)=\cos\!\Big(2\pi\frac{f+jf_s}{f_s}n+\varphi\Big)
=\cos\!\Big(2\pi\frac{f_s-f}{f_s}n-\varphi\Big).
\]
Consequently every sampled sinusoid coincides, sample for sample, with a sinusoid whose
frequency lies in $[0,f_s/2]$. The frequency $f_s/2$ is the \emph{Nyquist frequency}.
\end{proposition}
\begin{proof}
The middle expression adds $2\pi jn$ to the angle, a multiple of $2\pi$. For the
right one, $2\pi(f_s-f)n/f_s-\varphi=2\pi n-(2\pi fn/f_s+\varphi)$, and cosine is even
and $2\pi$-periodic. For the consequence, first subtract a multiple of $f_s$ to bring
$f$ into $[0,f_s)$; if it then exceeds $f_s/2$, replace it by $f_s-f\in(0,f_s/2)$.
\end{proof}

\subsection{Practice: by hand}

\begin{example}[The ocarina's five notes]
The ocarina plays D4, F4, A4, B4, D5, with MIDI numbers 62, 65, 69, 71, 74 and
frequencies $293.66$, $349.23$, $440.00$, $493.88$, $587.33$ Hz. A person humming
D4 an octave lower produces D3: $m=50$, $f=440\cdot2^{-19/12}=146.83$ Hz, half of
$293.66$.
\end{example}

\begin{example}[One frame]
At $f_s=44{,}100$ Hz, 40 ms is $0.04\times44{,}100=1{,}764$ samples. The lowest pitch
the code accepts is 80 Hz, whose period is $12.5$ ms, so a 40 ms frame holds $3.2$
periods: enough to see the repetition.
\end{example}

\begin{example}[An alias]
At $f_s=8{,}000$ Hz, a $5{,}000$ Hz cosine produces exactly the samples of a
$3{,}000$ Hz cosine (with its phase negated): $8{,}000-5{,}000=3{,}000\le4{,}000$.
After sampling, nothing can tell them apart.
\end{example}

\begin{connection}
\textbf{Code connection.} \code{songs.py} stores each song as a tuple of MIDI numbers
(Zelda's Lullaby is \code{(71, 74, 69, 71, 74, 69)}). \code{synth.render} builds
exactly the sum of Proposition~\ref{prop:harmperiod}, with harmonic amplitudes
$(1, 0.5, 0.25, 0.125)$ for a hum and $(1, 0.12, 0.04)$ for the nearly pure ocarina.
\code{DetectorConfig} searches pitches in $80$--$1{,}200$ Hz, far below the Nyquist
frequency $22{,}050$ Hz, and the sound card filters out content above $f_s/2$ before
sampling, so aliasing of the voice is not a concern.
\end{connection}

\subsection{Exercises}
\begin{exercise}[Hand]\label{ex:missing}
Let $x(t)=\cos(2\pi\cdot2f_0t)+\cos(2\pi\cdot3f_0t)$: harmonics 2 and 3 only, with no
component at $f_0$. Show that $1/f_0$ is a period, but $1/(2f_0)$ and $1/(3f_0)$ are
not. (This is the ``missing fundamental'' on slide 3: the pattern still repeats at
$f_0$, and the ear still hears $f_0$.)
\end{exercise}

\begin{exercise}[Hand]\label{ex:alias}
At $f_s=44{,}100$ Hz, which frequency in $[0,22{,}050]$ produces the same samples as a
$30{,}000$ Hz cosine? As a $50{,}000$ Hz cosine?
\end{exercise}

\begin{exercise}[Proof]\label{ex:nyqinj}
Show that if $0\le f,g\le f_s/2$ and $\cos(2\pi fn/f_s)=\cos(2\pi gn/f_s)$ for all
integers $n$, then $f=g$. So below the Nyquist frequency, sampled cosines are never
confused, and Proposition~\ref{prop:alias} is the whole story.
\end{exercise}

% =====================================================================
\section{Complex Exponentials and Probes}\label{sec:complex}
\usedin{every spectrum value in the code; the probe picture on slide 6.}

The central question of the project is: \emph{how much of frequency $f$ is in this
frame?} The answer is a dot product with a probe vector, and complex numbers are the
bookkeeping that makes the answer independent of phase.

\subsection{Theory}

Recall Euler's formula $e^{i\theta}=\cos\theta+i\sin\theta$: the point at angle
$\theta$ on the unit circle. Hence $\abs{e^{i\theta}}=1$,
$\overline{e^{i\theta}}=e^{-i\theta}$, $e^{i\theta}e^{i\psi}=e^{i(\theta+\psi)}$, and
\[
\cos\theta=\frac{e^{i\theta}+e^{-i\theta}}{2},\qquad
\sin\theta=\frac{e^{i\theta}-e^{-i\theta}}{2i}.
\]

\begin{definition}[Inner product on $\C^N$]
For $x,y\in\C^N$, $\ip{x}{y}=\sum_{n}x_n\overline{y_n}$ and $\norm{x}=\sqrt{\ip xx}$.
\end{definition}

The conjugate makes $\ip xx=\sum\abs{x_n}^2\ge0$ a genuine squared length. The price:
$\ip yx=\overline{\ip xy}$, and $\ip{x}{\lambda y}=\bar\lambda\ip xy$ while
$\ip{\lambda x}{y}=\lambda\ip xy$. On real vectors it is the ordinary dot product.

\begin{lemma}[Roots of unity]\label{lem:roots}
Let $N\ge1$ and $\omega=e^{-2\pi i/N}$. For every integer $m$,
\[
\sum_{n=0}^{N-1}\omega^{mn}=\begin{cases}N&\text{if }N\mid m,\\0&\text{otherwise.}\end{cases}
\]
\end{lemma}
\begin{proof}
If $N\mid m$ then $\omega^m=e^{-2\pi i m/N}=1$ and all $N$ terms are 1. Otherwise
$r=\omega^m\ne1$ but $r^N=e^{-2\pi i m}=1$, so the geometric series gives
$\sum_{n=0}^{N-1}r^n=(1-r^N)/(1-r)=0$.
\end{proof}

Geometrically, the terms are equally spaced points on the unit circle, and they
balance. Every orthogonality fact in this booklet is this lemma.

\begin{proposition}[A complex probe is a cosine probe and a sine probe]\label{prop:probe}
For real $x\in\R^N$ and real angles $\theta_0,\dots,\theta_{N-1}$,
\[
\Big\lvert\sum_{n}x_ne^{-i\theta_n}\Big\rvert^2
=\Big(\sum_nx_n\cos\theta_n\Big)^2+\Big(\sum_nx_n\sin\theta_n\Big)^2.
\]
\end{proposition}
\begin{proof}
$e^{-i\theta_n}=\cos\theta_n-i\sin\theta_n$, and $x_n$ is real, so the sum has real
part $\sum x_n\cos\theta_n$ and imaginary part $-\sum x_n\sin\theta_n$. Apply
$\abs{a+bi}^2=a^2+b^2$.
\end{proof}

With $\theta_n=2\pi fn/f_s$, the left side is the \emph{score} of frequency $f$ on
slide 6: a cosine probe alone would miss a sine-shaped component (their products
cancel), but the pair cannot miss it. For probe frequencies on the grid $f=kf_s/N$,
the score is exactly phase-blind:

\begin{proposition}[Score of a sampled cosine]\label{prop:cosbin}
Let $0<k<N/2$ be an integer and $x_n=a\cos(2\pi kn/N+\varphi)$. Then
\[
\sum_{n=0}^{N-1}x_n\,\omega^{kn}=\frac{aN}{2}e^{i\varphi},
\qquad\text{so the score is }\frac{aN}{2}\text{ whatever }\varphi\text{ is.}
\]
\end{proposition}
\begin{proof}
Write $x_n=\frac a2\big(e^{i\varphi}\omega^{-kn}+e^{-i\varphi}\omega^{kn}\big)$. Then
$\sum_nx_n\omega^{kn}=\frac a2\big(e^{i\varphi}\sum_n1+e^{-i\varphi}\sum_n\omega^{2kn}\big)$.
The first sum is $N$. Since $0<2k<N$, $N\nmid 2k$, and Lemma~\ref{lem:roots} makes the
second sum 0.
\end{proof}

\subsection{Practice: by hand}

\begin{example}[Fourth roots of unity]
For $N=4$, $\omega=e^{-2\pi i/4}=-i$, with powers $1,-i,-1,i$: the four compass
points. For $m=1$: $1-i-1+i=0$. For $m=4$: $1+1+1+1=4$.
\end{example}

\begin{example}[Why two probes]
Let $N=4$ and probe one cycle per frame, $\theta_n=2\pi n/4$. The cosine
$x=(1,0,-1,0)$ has cosine score $1\cdot1+(-1)(-1)=2$ and sine score 0. The sine
$y=(0,1,0,-1)$ has cosine score \emph{0}: a cosine probe alone would say ``no energy
here''. Its sine score is $1+1=2$. The complex probe gives
$\abs{\sum y_n(-i)^n}=\abs{-i-i}=2$ for both, as
Proposition~\ref{prop:cosbin} promises with $a=1$, $N=4$.
\end{example}

\begin{connection}
\textbf{Code connection.} The code never loops over probe frequencies. It computes all
the grid scores at once with an FFT (Lessons~\ref{sec:dft} and~\ref{sec:fft}), and
\code{magnitude\_spectrum} keeps only their magnitudes, which by
Proposition~\ref{prop:cosbin} do not depend on phase. One probe by hand would be
\code{abs(x @ np.exp(-2j * np.pi * f * np.arange(N) / fs))}.
\end{connection}

\subsection{Exercises}
\begin{exercise}[Hand]\label{ex:sixth}
List the sixth roots of unity $e^{-2\pi in/6}$, $n=0,\dots,5$, as $a+bi$. Check that
they sum to 0 and that their squares also sum to 0. Which case of
Lemma~\ref{lem:roots} is each?
\end{exercise}

\begin{exercise}[Hand]\label{ex:probe4}
For $N=4$, compute $\sum_nx_n(-i)^n$ for $x=(1,1,-1,-1)$. Write $x_n$ as
$a\cos(2\pi n/4+\varphi)$ and check your answer against
Proposition~\ref{prop:cosbin}.
\end{exercise}

\begin{exercise}[Proof, $\star$]\label{ex:edgebins}
Redo Proposition~\ref{prop:cosbin} for $k=0$ and, when $N$ is even, for $k=N/2$. Show
that the sum is $aN\cos\varphi$ in both cases, so at these two bins the score
\emph{does} depend on the phase. What is special about $\cos(\pi n+\varphi)$?
\end{exercise}

% =====================================================================
\section{The DFT Is a Change of Basis}\label{sec:dft}
\usedin{\code{magnitude\_spectrum} (\code{fft.rfft}), and \code{np.fft} generally.
Slide 7.}

Take $N$ probes on a grid, one for each whole number of cycles per frame, and ask all
$N$ questions at once. The answer is a matrix--vector product, and because the probes
are orthogonal, nothing is lost: it is a change of basis.

\subsection{Theory}

\begin{definition}[Fourier vectors, DFT]
For $k=0,\dots,N-1$, the \emph{Fourier vector} $e_k\in\C^N$ has entries
$(e_k)_n=e^{2\pi ikn/N}$: it turns $k$ times around the unit circle across the
frame. The \emph{discrete Fourier transform} of $x\in\C^N$ is $X\in\C^N$ with
\[
X_k=\ip{x}{e_k}=\sum_{n=0}^{N-1}x_n\,\omega^{kn},\qquad\omega=e^{-2\pi i/N}.
\]
Each $X_k$ is called a \emph{bin}. In matrix form $X=Fx$ with $F_{kn}=\omega^{kn}$; the
rows of $F$ are the conjugates $\overline{e_k}^{\,\mathsf T}$.
\end{definition}

\begin{theorem}[Orthogonality of the Fourier vectors]\label{thm:orth}
For $0\le j,k\le N-1$: $\ip{e_j}{e_k}=N$ if $j=k$, and $0$ otherwise.
\end{theorem}
\begin{proof}
$\ip{e_j}{e_k}=\sum_ne^{2\pi ijn/N}e^{-2\pi ikn/N}=\sum_n\omega^{(k-j)n}$. By
Lemma~\ref{lem:roots} this is $N$ if $N\mid(k-j)$ and $0$ otherwise. Since
$\abs{k-j}<N$, $N\mid(k-j)$ only when $j=k$.
\end{proof}

\begin{corollary}[Unitary up to $N$; inversion; Parseval]\label{cor:unitary}
$F^*F=FF^*=N\,I$, where $F^*=\overline F^{\mathsf T}$. Hence $F/\sqrt N$ is unitary,
$F^{-1}=F^*/N$, and
\[
x=\frac1N\sum_{k=0}^{N-1}X_k\,e_k,\qquad\norm{X}^2=N\norm{x}^2.
\]
\end{corollary}
\begin{proof}
$(FF^*)_{jk}=\sum_nF_{jn}\overline{F_{kn}}=\sum_n\overline{(e_j)_n}(e_k)_n
=\ip{e_k}{e_j}$, which is $N\delta_{jk}$ by Theorem~\ref{thm:orth}. So $FF^*=NI$; for
square matrices a one-sided inverse is two-sided, so $F^*F=NI$ too. Then
$F^{-1}=F^*/N$, and $x=F^{-1}X=\frac1NF^*X$ is the displayed sum because the columns
of $F^*$ are the $e_k$. Finally $\norm{Fx}^2=(Fx)^*(Fx)=x^*F^*Fx=N\norm x^2$.
\end{proof}

So $X$ is just $x$ in new coordinates: the coefficients of $x$ along the orthogonal
basis $e_0,\dots,e_{N-1}$, up to the factor $N$. This is the \emph{frequency domain};
the \emph{Fourier transform} is the change of coordinates between the two domains.

\begin{proposition}[Bins are frequencies]
If $x_n=x(n/f_s)$, then $(e_k)_n=e^{2\pi i(kf_s/N)(n/f_s)}$ is the complex sinusoid of
frequency $kf_s/N$ Hz, sampled at the same instants. So bin $k$ measures frequency
$kf_s/N$, and neighbouring bins are $f_s/N$ Hz apart.
\end{proposition}
\begin{proof}
$2\pi i(kf_s/N)(n/f_s)=2\pi ikn/N$.
\end{proof}

\subsection{Practice: by hand}

\begin{example}[The $4\times4$ DFT]
With $\omega=-i$,
\[
F=\begin{pmatrix}1&1&1&1\\1&-i&-1&i\\1&-1&1&-1\\1&i&-1&-i\end{pmatrix},\qquad
F\begin{pmatrix}1\\2\\3\\4\end{pmatrix}=\begin{pmatrix}10\\-2+2i\\-2\\-2-2i\end{pmatrix}.
\]
Parseval: $\norm x^2=1+4+9+16=30$ and $\norm X^2=100+8+4+8=120=4\cdot30$. Inversion at
$n=0$: $\frac14(10+(-2+2i)+(-2)+(-2-2i))=\frac44=1=x_0$.
\end{example}

\begin{example}[Two extremes]
The impulse $(1,0,0,0)$ has DFT $(1,1,1,1)$: one instant contains every frequency
equally. The constant $(1,1,1,1)$ has DFT $(4,0,0,0)$: all of it is in bin 0.
\end{example}

\begin{connection}
\textbf{Code connection.} \code{np.fft.fft(x)} computes $Fx$; \code{np.fft.ifft(X)}
computes $F^*X/N$, so NumPy puts the whole $1/N$ in the inverse;
\code{np.fft.fft(x, norm="ortho")} computes the unitary $(F/\sqrt N)x$. In
\code{magnitude\_spectrum} the transform length is $7{,}200$ at $f_s=44{,}100$, so
bins are $44{,}100/7{,}200=6.125$ Hz apart.
\end{connection}

\subsection{Exercises}
\begin{exercise}[Hand]\label{ex:dft4}
Compute the DFT of $(1,0,-1,0)$ and of $(0,1,0,-1)$, and compare their magnitudes using
Lesson~\ref{sec:complex}. Then invert $X=(0,2,0,2)$ with the formula of
Corollary~\ref{cor:unitary} and confirm you get the first signal back.
\end{exercise}

\begin{exercise}[Proof]\label{ex:parseval}
Prove $\ip{Fx}{Fy}=N\ip xy$ for all $x,y\in\C^N$, and deduce Parseval.
\end{exercise}

\begin{exercise}[Proof, $\star$]\label{ex:basis}
Without using matrices, show that $e_0,\dots,e_{N-1}$ is a basis of $\C^N$: prove that
pairwise orthogonal nonzero vectors are linearly independent (take the inner product of
a vanishing combination with each $e_j$), then count.
\end{exercise}

% =====================================================================
\section{Reading the DFT: Periodicity, Real Signals, \texttt{rfft}}\label{sec:layout}
\usedin{\code{magnitude\_spectrum} (\code{fft.rfft} returns $3{,}601$ bins);
\code{harmonic\_product} (the bound \code{n\_bins}). Slides 8 and 10.}

The array \code{np.fft.fft} returns has a particular layout: zero frequency first,
then positive frequencies, then \emph{negative} ones. This lesson explains every part
of that layout from one fact: the DFT treats your $N$ samples as one period of an
endless loop.

\subsection{Theory}

\begin{theorem}[The DFT sees a loop]\label{thm:periodic}
Define $X_k=\sum_{n=0}^{N-1}x_n\omega^{kn}$ for \emph{every} integer $k$, and
$\tilde x_n=\frac1N\sum_{k=0}^{N-1}X_k e^{2\pi ikn/N}$ for every integer $n$. Then
\[
X_{k+N}=X_k\quad\text{for all }k,\qquad
\tilde x_n=x_n\ (0\le n<N),\qquad \tilde x_{n+N}=\tilde x_n\quad\text{for all }n.
\]
\end{theorem}
\begin{proof}
$\omega^{(k+N)n}=\omega^{kn}\omega^{Nn}=\omega^{kn}$ because $\omega^N=1$. The same
argument applies to $e^{2\pi i(n+N)k/N}=e^{2\pi ikn/N}$. And $\tilde x_n=x_n$ for
$0\le n<N$ is Corollary~\ref{cor:unitary}.
\end{proof}

The sequence $\tilde x$ is the \emph{periodic extension} of $x$: the frame tiled end to
end. Any finite list can be tiled this way, so the DFT can always be applied, but it
silently analyses the tiled signal, seams included (Lesson~\ref{sec:window}).

Periodicity in $k$ is aliasing (Proposition~\ref{prop:alias}) in disguise: frequencies
$k$ and $k+N$ cycles per frame give the same samples. In particular $X_{N-k}=X_{-k}$,
so the upper half of the array holds the \emph{negative frequencies}.

\begin{definition}[DC, positive, Nyquist, negative bins]
Bin 0 is the \emph{zero frequency} or \emph{DC} bin: $X_0=\sum x_n$, which is $N$ times
the mean. Bins $1,\dots,\lceil N/2\rceil-1$ are positive frequencies. If $N$ is even,
bin $N/2$ is the \emph{Nyquist} bin, frequency $f_s/2$, whose Fourier vector
$(e_{N/2})_n=(-1)^n$ is the fastest real oscillation the samples can hold. The
remaining bins $k$ stand for the negative frequencies $k-N$. Odd $N$ has no Nyquist
bin.
\end{definition}

\begin{theorem}[Real signals have conjugate-symmetric spectra]\label{thm:realsym}
If $x\in\R^N$, then $X_{N-k}=\overline{X_k}$ for every $k$. In particular $X_0$ is real,
and when $N$ is even, $X_{N/2}$ is real.
\end{theorem}
\begin{proof}
$X_{N-k}=X_{-k}=\sum_nx_n\omega^{-kn}=\sum_n\overline{x_n\omega^{kn}}=\overline{X_k}$,
using that $x_n$ is real and $\overline{\omega}=\omega^{-1}$. For $k=0$ this says
$X_0=\overline{X_0}$; for $k=N/2$ it says $X_{N/2}=\overline{X_{N/2}}$.
\end{proof}

\begin{corollary}[Why \code{rfft} returns $\lfloor N/2\rfloor+1$ bins]\label{cor:rfft}
For real $x$, the bins $X_0,\dots,X_{\lfloor N/2\rfloor}$ determine all of $X$, hence
$x$. Counting real parameters: $X_0$ (1 real number), $X_{N/2}$ if $N$ is even
(1 real number), and $2$ real numbers for each remaining bin, which totals exactly $N$.
\end{corollary}
\begin{proof}
The other bins are conjugates of these by Theorem~\ref{thm:realsym}, and $x=F^{-1}X$.
For even $N$: $1+1+2(N/2-1)=N$. For odd $N$: $1+2\cdot\frac{N-1}2=N$.
\end{proof}

\subsection{Practice: by hand}

\begin{example}[The layout for $N=8$]
Read the eight bins of \code{np.fft.fft} like this:
\begin{center}
\small
\begin{tabular}{@{}lcccccccc@{}}
\toprule
bin $k$ & 0 & 1 & 2 & 3 & 4 & 5 & 6 & 7\\
cycles per frame & 0 & $+1$ & $+2$ & $+3$ & $\pm4$ & $-3$ & $-2$ & $-1$\\
meaning & DC & pos & pos & pos & Nyquist & neg & neg & neg\\
\bottomrule
\end{tabular}
\end{center}
\code{np.fft.fftfreq(8)} labels bin 4 as $-0.5$ (cycles per sample): by periodicity,
$+4$ and $-4$ are the same bin. \code{np.fft.rfft} keeps bins 0--4.
\end{example}

\begin{example}[Symmetry in action]
The DFT of $(1,2,3,4)$ was $(10,\,-2+2i,\,-2,\,-2-2i)$: $X_3=\overline{X_1}$, and
$X_0$, $X_2$ are real.
\end{example}

\begin{example}[The project's numbers]
Frames are zero-padded to $N=7{,}200$ (Lesson~\ref{sec:window}), so \code{rfft}
returns $7{,}200/2+1=3{,}601$ bins, at $0,6.125,\dots,22{,}050$ Hz. A hummed D3
($146.83$ Hz) lands near bin $146.83/6.125\approx23.97$.
\end{example}

\begin{connection}
\textbf{Code connection.} \code{spectrum = np.abs(fft.rfft(frames * np.hanning(n),
n=n\_fft, axis=-1))} keeps only the non-redundant half (Corollary~\ref{cor:rfft}),
and the magnitudes $\abs{X_k}$, discarding phase. \code{harmonic\_product} uses
$\code{n\_bins} = 3{,}601 /\!/ 3 = 1{,}200$, so that bin $3k$ exists for every
candidate $k<1{,}200$.
\end{connection}

\subsection{Exercises}
\begin{exercise}[Hand]\label{ex:layout7}
(a) Make the layout table above for $N=7$. Which bin is the Nyquist bin, and how many
bins does \code{rfft} return? (b) A real signal of length 4 has $X_0=10$, $X_1=-2+2i$,
$X_2=-2$. Find $X_3$, then $x$.
\end{exercise}

\begin{exercise}[Hand]\label{ex:nyq}
For $N=4$, compute the DFT of $x_n=\cos(\pi n+\varphi)$ for $\varphi=0$ and
$\varphi=\pi/2$. Explain why the Nyquist bin cannot record phase (compare
Exercise~\ref{ex:edgebins}).
\end{exercise}

\begin{exercise}[Proof, $\star$]\label{ex:realseries}
Show that every real $x\in\R^N$, $N$ even, can be written as
\[
x_n=\frac{X_0}N+\sum_{k=1}^{N/2-1}\frac{2\abs{X_k}}N\cos\!\Big(\frac{2\pi kn}N+\arg X_k\Big)
+\frac{X_{N/2}}N(-1)^n.
\]
This is the real-world meaning of the bins: amplitude $2\abs{X_k}/N$ and phase
$\arg X_k$.
\end{exercise}

% =====================================================================
\section{Shifts, Circulants, and Convolution}\label{sec:circ}
\usedin{the decision to keep only $\abs{X}$; the Hann window of
Lesson~\ref{sec:window}. Slide 9.}

Why sinusoids, and not some other orthogonal basis? Because sinusoids are the
eigenvectors of \emph{shifting}. That one fact explains why the magnitude spectrum
does not care where in the frame a note sits, and why multiplying by a window blurs
the spectrum.

\subsection{Theory}

\begin{definition}[Cyclic shift, circulant, cyclic convolution]
The \emph{cyclic shift} $S$ on $\C^N$ is $(Sx)_n=x_{n-1}$, indices mod $N$. A
\emph{circulant} matrix $C$ has first column $c$ and each later column equal to the
previous one shifted down cyclically: $C_{nj}=c_{n-j}$. The \emph{cyclic convolution}
of $c$ and $x$ is $(c\cconv x)_n=\sum_{j=0}^{N-1}c_jx_{n-j}$.
\end{definition}

Note that $Cx=c\cconv x$ directly from $C_{nj}=c_{n-j}$ (substitute $j\mapsto n-j$),
and $C=\sum_jc_jS^j$, since $S^j$ is the circulant with first column the $j$-th
standard basis vector.

\begin{lemma}[Fourier vectors are eigenvectors of the shift]\label{lem:shifteig}
$Se_k=\omega^ke_k$ for every $k$.
\end{lemma}
\begin{proof}
$(Se_k)_n=(e_k)_{n-1}=e^{2\pi ik(n-1)/N}=e^{-2\pi ik/N}e^{2\pi ikn/N}=\omega^k(e_k)_n$.
(Indices mod $N$ are harmless because $e_k$ has period $N$.)
\end{proof}

\begin{theorem}[Shift theorem]\label{thm:shift}
$(FS^mx)_k=\omega^{km}X_k$. Consequently $\abs{(FS^mx)_k}=\abs{X_k}$: a cyclic shift
changes only phases, never magnitudes.
\end{theorem}
\begin{proof}
$(FS^mx)_k=\ip{S^mx}{e_k}=\ip{x}{(S^*)^me_k}$. The shift is unitary with
$S^*=S^{-1}$, and Lemma~\ref{lem:shifteig} gives $S^{-1}e_k=\omega^{-k}e_k$, so
$(S^*)^me_k=\omega^{-km}e_k$. Conjugate-linearity in the second slot turns
$\omega^{-km}$ into $\overline{\omega^{-km}}=\omega^{km}$, giving
$\omega^{km}\ip{x}{e_k}$. Finally $\abs{\omega^{km}}=1$.
\end{proof}

\begin{theorem}[The DFT diagonalizes every circulant]\label{thm:diag}
For the circulant $C$ with first column $c$, $Ce_k=(Fc)_k\,e_k$. Equivalently
$FCF^{-1}=\diag(Fc)$, and
\[
F(c\cconv x)=(Fc)\odot(Fx)\qquad\text{(the convolution theorem),}
\]
where $\odot$ is the entrywise product.
\end{theorem}
\begin{proof}
$Ce_k=\sum_jc_jS^je_k=\sum_jc_j\omega^{jk}e_k=(Fc)_ke_k$ by Lemma~\ref{lem:shifteig}.
So the columns of $F^{-1}=F^*/N$, which are multiples of the $e_k$, are eigenvectors
of $C$ with eigenvalues $(Fc)_k$; that is $CF^{-1}=F^{-1}\diag(Fc)$. Multiply on the
left by $F$, then apply both sides to $Fx$:
$F(c\cconv x)=FCx=\diag(Fc)Fx$.
\end{proof}

The dual statement swaps the roles of the two domains.

\begin{theorem}[Products become convolutions]\label{thm:dual}
For $x,y\in\C^N$, $F(x\odot y)=\frac1N(Fx)\cconv(Fy)$.
\end{theorem}
\begin{proof}
Write $y_n=\frac1N\sum_jY_je^{2\pi ijn/N}$. Then
\[
(F(x\odot y))_k=\sum_nx_n\frac1N\sum_jY_je^{2\pi ijn/N}e^{-2\pi ikn/N}
=\frac1N\sum_jY_j\sum_nx_n\omega^{(k-j)n}=\frac1N\sum_jY_jX_{k-j},
\]
using Theorem~\ref{thm:periodic} for the index $k-j$.
\end{proof}

\subsection{Practice: by hand}

\begin{example}[Convolution theorem with $N=4$]
Take $c=(1,1,0,0)$, so $(c\cconv x)_n=x_n+x_{n-1}$, and $x=(1,2,3,4)$. Then
$c\cconv x=(5,3,5,7)$. On the other side, $Fc=(2,1-i,0,1+i)$ and
$Fx=(10,-2+2i,-2,-2-2i)$, whose entrywise product is $(20,4i,0,-4i)$. Directly,
$F(5,3,5,7)=(20,4i,0,-4i)$.
\end{example}

\begin{example}[A blur and its eigenvalues]
The circulant with $c=(\frac12,\frac14,0,\frac14)$ replaces each entry by a weighted
average with its neighbours. Its eigenvalues are
$(Fc)_k=\frac12+\frac14\omega^k+\frac14\omega^{-k}=\frac12+\frac12\cos(2\pi k/4)$,
that is $(1,\frac12,0,\frac12)$: slow oscillations pass, the fastest is erased. A blur
is a low-pass filter.
\end{example}

\begin{connection}
\textbf{Code connection.} Consecutive frames are 10 ms apart, so the same note sits at
different positions in different frames. Theorem~\ref{thm:shift} is why
\code{magnitude\_spectrum} keeps $\abs{X}$: position becomes phase, and phase is
thrown away. (A real shift is not exactly cyclic, since samples enter and leave the
frame, so this holds approximately; the window of Lesson~\ref{sec:window} keeps the
approximation good.)
\end{connection}

\subsection{Exercises}
\begin{exercise}[Hand]\label{ex:diff}
Find the eigenvalues of the circulant with $c=(1,-1,0,0)$, that is $x_n-x_{n-1}$.
Which frequency is erased, and which is amplified most?
\end{exercise}

\begin{exercise}[Proof]\label{ex:sunitary}
Show that $S$ is unitary and $S^N=I$, and use Lemma~\ref{lem:shifteig} to deduce that
its eigenvalues are exactly the $N$-th roots of unity.
\end{exercise}

\begin{exercise}[Proof, $\star$]\label{ex:commute}
Show that a matrix $A$ commutes with $S$ if and only if $A$ is circulant. So
``shift-invariant linear map'' and ``circulant'' are the same thing, and
Theorem~\ref{thm:diag} diagonalizes all of them at once.
\end{exercise}

% =====================================================================
\section{Leakage, Windows, and Zero-Padding}\label{sec:window}
\usedin{\code{magnitude\_spectrum}: \code{frames * np.hanning(n)} and
\code{n=n\_fft} with \code{zero\_pad=4}. Slides 11 and 14.}

A hummed note almost never completes a whole number of cycles in a frame. Its
periodic extension then has a jump at every seam, and its energy spreads over every
bin. This lesson measures the spread exactly, shows how a window tames it, and shows
what zero-padding does and does not do.

\subsection{Theory}

\begin{theorem}[Leakage]\label{thm:leak}
Let $x_n=e^{2\pi i\alpha n/N}$ with $\alpha\in\R$. If $\alpha$ is an integer, $X$ is
$N$ at bin $\alpha\bmod N$ and $0$ elsewhere. Otherwise, for every $k$,
\[
\abs{X_k}=\frac{\abs{\sin\pi(\alpha-k)}}{\abs{\sin\big(\pi(\alpha-k)/N\big)}}\ne0.
\]
\end{theorem}
\begin{proof}
$X_k=\sum_n r^n$ with $r=e^{2\pi i(\alpha-k)/N}$. If $\alpha$ is an integer this is
Lemma~\ref{lem:roots}. Otherwise $r\ne1$, and
$X_k=(1-r^N)/(1-r)=(1-e^{2\pi i(\alpha-k)})/(1-e^{2\pi i(\alpha-k)/N})$. Use
$\abs{1-e^{i\theta}}=2\abs{\sin(\theta/2)}$ on the top and the bottom.
\end{proof}

For $\abs{\alpha-k}$ small compared with $N$, the denominator is about
$\pi\abs{\alpha-k}/N$, so $\abs{X_k}\approx N\abs{\sin\pi(\alpha-k)}/(\pi\abs{\alpha-k})$:
the leaked energy decays only like $1/\text{distance}$. That slow decay is what a
window fixes.

\begin{definition}[Window, Hann window]
A \emph{window} is a vector $w$ that multiplies the frame entrywise before the DFT:
$x\mapsto w\odot x=\diag(w)x$. The (periodic) \emph{Hann window} is
$w_n=\frac12-\frac12\cos(2\pi n/N)$, a smooth bump that is 0 at $n=0$ and fades both
ends of the loop, removing the jump at the seam.
\end{definition}

\begin{proposition}[The Hann window has three nonzero bins]\label{prop:hann}
$\hat w=Fw$ is $N/2$ at bin 0, $-N/4$ at bins $1$ and $N-1$, and 0 elsewhere.
\end{proposition}
\begin{proof}
$w=\frac12e_0-\frac14e_1-\frac14e_{N-1}$, since
$\cos(2\pi n/N)=\frac12(e^{2\pi in/N}+e^{2\pi i(N-1)n/N})$. Apply $F$ and
Theorem~\ref{thm:orth}: $Fe_j=N\cdot(\text{$j$-th standard basis vector})$.
\end{proof}

\begin{theorem}[Windowing is a three-tap blur of the spectrum]\label{thm:hannblur}
With the Hann window,
\[
\big(F(w\odot x)\big)_k=\tfrac12X_k-\tfrac14X_{k-1}-\tfrac14X_{k+1}.
\]
\end{theorem}
\begin{proof}
By Theorem~\ref{thm:dual}, $F(w\odot x)=\frac1N\hat w\cconv X$, and by
Proposition~\ref{prop:hann} only three terms of the convolution survive:
$\frac1N(\frac N2X_k-\frac N4X_{k-1}-\frac N4X_{k+1})$.
\end{proof}

Why this helps: far from the peak, $1-e^{2\pi i(\alpha-k)/N}\approx-2\pi i(\alpha-k)/N$,
so the leaked bins behave like $X_k\approx K/(\alpha-k)$ for a constant $K$, a slowly
varying function of $k$. The three taps are a negative second difference,
$\frac12X_k-\frac14(X_{k-1}+X_{k+1})=-\frac14\big(X_{k+1}-2X_k+X_{k-1}\big)$, and the
second difference of $1/u$ is about $2/u^3$. So the tails fall like $1/\text{distance}^3$
instead of $1/\text{distance}$: the highest sidelobe drops from about $-13$ dB to $-31$ dB
below the peak (a decibel is $20\log_{10}$ of an amplitude ratio). The price is a peak
twice as wide.

\begin{definition}[Spectrum of a finite sequence]
For $x\in\C^N$ and real $\nu$ (cycles per sample), let
$\mathcal X(\nu)=\sum_{n=0}^{N-1}x_ne^{-2\pi i\nu n}$. This is a smooth, 1-periodic
function of $\nu$ (its formal name is the discrete-time Fourier transform), and
$X_k=\mathcal X(k/N)$.
\end{definition}

\begin{theorem}[Zero-padding samples the same curve more finely]\label{thm:zeropad}
Append zeros to $x$ to length $M\ge N$ and take the $M$-point DFT. Its bins are
$\mathcal X(k/M)$, $k=0,\dots,M-1$.
\end{theorem}
\begin{proof}
The padded entries are 0, so $\sum_{n=0}^{M-1}x_ne^{-2\pi ikn/M}
=\sum_{n=0}^{N-1}x_ne^{-2\pi i(k/M)n}=\mathcal X(k/M)$.
\end{proof}

So zero-padding adds no information; it evaluates the same curve $\mathcal X$ on a grid
$M/N$ times finer. That is exactly what peak-finding needs (Lesson~\ref{sec:pitch}).

\subsection{Practice: by hand}

\begin{example}[Half a cycle too many]
$N=4$, $\alpha=\frac12$. Theorem~\ref{thm:leak} gives
\[
\abs{X}=\Big(\frac1{\sin\frac\pi8},\frac1{\sin\frac\pi8},\frac1{\sin\frac{3\pi}8},
\frac1{\sin\frac{3\pi}8}\Big)\approx(2.613,\,2.613,\,1.082,\,1.082).
\]
Every bin is nonzero. Parseval check: $2(6.828)+2(1.172)=16=N\norm x^2$.
\end{example}

\begin{example}[Hann on a bin-centred cosine]
$N=8$ and $x_n=\cos(2\pi\cdot2n/8)$, so $X=(0,0,4,0,0,0,4,0)$. By
Theorem~\ref{thm:hannblur},
\[
F(w\odot x)=(0,-1,2,-1,0,-1,2,-1):
\]
each spike becomes the three taps $-\frac14,\frac12,-\frac14$, times 4.
\end{example}

\begin{connection}
\textbf{Code connection.} \code{magnitude\_spectrum} multiplies each 1{,}764-sample
frame by \code{np.hanning(n)}, the \emph{symmetric} Hann window
$\frac12-\frac12\cos(2\pi n/(N-1))$; at $N=1{,}764$ it is indistinguishable from the
periodic one analysed here. It then calls \code{rfft(..., n=n\_fft)} with
$n\_fft=7{,}200$ (from \code{next\_fast\_len(4 * 1764)}), which zero-pads by
Theorem~\ref{thm:zeropad}: the grid goes from 25 Hz to 6.125 Hz.
\end{connection}

\subsection{Exercises}
\begin{exercise}[Hand]\label{ex:hann4}
Hann-window $x=(1,2,3,4)$ with $N=4$, so $w=(0,\frac12,1,\frac12)$. Compute
$F(w\odot x)$ directly and check it against Theorem~\ref{thm:hannblur} using
$X=(10,-2+2i,-2,-2-2i)$.
\end{exercise}

\begin{exercise}[Hand]\label{ex:leak8}
For $N=8$ and $\alpha=2.5$, compute $\abs{X_2}$, $\abs{X_3}$ and $\abs{X_6}$ from
Theorem~\ref{thm:leak}.
\end{exercise}

\begin{exercise}[Proof, $\star$]\label{ex:decay}
Using $\sin y\ge 2y/\pi$ for $0\le y\le\pi/2$, show that
$\abs{X_k}\le N/(2\abs{\alpha-k})$ whenever $0<\abs{\alpha-k}\le N/2$. This is the
``decays like $1/\text{distance}$'' claim made precise.
\end{exercise}

% =====================================================================
\section{The Fast Fourier Transform}\label{sec:fft}
\usedin{\code{scipy.fft.rfft} inside \code{magnitude\_spectrum};
\code{next\_fast\_len}. Slide 12.}

$Fx$ costs $N^2$ multiply--adds done directly: about 52 million for $N=7{,}200$, a
hundred times a second. The FFT computes the same vector in about $N\log_2N$
operations by divide and conquer, exactly like merge sort.

\subsection{Theory}

\begin{theorem}[Radix-2 split]\label{thm:radix2}
Let $N$ be even, and let $E$ and $O$ be the $N/2$-point DFTs of the even-indexed
samples $(x_0,x_2,\dots)$ and the odd-indexed samples $(x_1,x_3,\dots)$. With
$\omega_N=e^{-2\pi i/N}$, for $0\le k<N/2$:
\[
X_k=E_k+\omega_N^kO_k,\qquad X_{k+N/2}=E_k-\omega_N^kO_k.
\]
\end{theorem}
\begin{proof}
Split the sum by parity, $n=2r$ or $n=2r+1$:
\[
X_k=\sum_{r}x_{2r}\omega_N^{2rk}+\omega_N^k\sum_rx_{2r+1}\omega_N^{2rk}
=E_k+\omega_N^kO_k,
\]
since $\omega_N^2=\omega_{N/2}$. Replacing $k$ by $k+N/2$: $E$ and $O$ have period
$N/2$ (Theorem~\ref{thm:periodic}), and $\omega_N^{N/2}=e^{-\pi i}=-1$.
\end{proof}

Each pair $(E_k,O_k)\mapsto(X_k,X_{k+N/2})$ is a \emph{butterfly}: one multiplication
by $\omega_N^k$, one addition, one subtraction.

\begin{theorem}[Cost]\label{thm:fftcost}
Let $T(N)$ count the additions and subtractions, so $T(1)=0$ and $T(N)=2T(N/2)+N$
(two half-size transforms, then $N/2$ butterflies with one addition and one subtraction
each). Then $T(2^p)=p\,2^p$, that is $T(N)=N\log_2N$. The multiplications by
$\omega_N^k$ add $\frac N2\log_2N$ more.
\end{theorem}
\begin{proof}
Induction on $p$. $T(1)=0=0\cdot1$. If $T(2^{p-1})=(p-1)2^{p-1}$ then
$T(2^p)=2(p-1)2^{p-1}+2^p=p\,2^p$.
\end{proof}

In matrix language the split says
$F_N=\begin{psmallmatrix}I&\Omega\\I&-\Omega\end{psmallmatrix}
\begin{psmallmatrix}F_{N/2}&0\\0&F_{N/2}\end{psmallmatrix}P$, with $P$ the
even/odd permutation and $\Omega=\diag(\omega_N^k)$: recursing factors $F$ into
$\log_2N$ sparse matrices with two nonzeros per row.

\subsection{Practice: by hand}

\begin{example}[A 4-point FFT]
For $x=(1,2,3,4)$: $E=\mathrm{DFT}_2(1,3)=(4,-2)$, $O=\mathrm{DFT}_2(2,4)=(6,-2)$,
$\omega_4=-i$. Then $X_0=4+6=10$, $X_1=-2+(-i)(-2)=-2+2i$, $X_2=4-6=-2$,
$X_3=-2-(-i)(-2)=-2-2i$: the same as Lesson~\ref{sec:dft}, with 8 additions and 2 multiplications
instead of 16 multiply--adds.
\end{example}

\begin{connection}
\textbf{Code connection.} The frames are padded to $7{,}200=2^5\cdot3^2\cdot5^2$, not a
power of two: \code{next\_fast\_len(7056, real=True)} picks the smallest length
$\ge7{,}056$ made only of 2s, 3s and 5s, because SciPy's FFT splits by 3 and 5 the same
way it splits by 2 (Exercise~\ref{ex:radix3}). Cost per frame: about
$7{,}200\log_27{,}200\approx92{,}000$ operations instead of $51{,}840{,}000$, about
$560\times$ fewer.
\end{connection}

\subsection{Exercises}
\begin{exercise}[Hand]\label{ex:fft4}
(a) Run the 4-point FFT on $(1,0,-1,0)$ and compare with Lesson~\ref{sec:dft}.
(b) Recursing the even/odd split three times for $N=8$ reorders the inputs. Show the
order is $x_0,x_4,x_2,x_6,x_1,x_5,x_3,x_7$, and explain the name \emph{bit reversal}.
\end{exercise}

\begin{exercise}[Proof, $\star$]\label{ex:radix3}
Let $3\mid N$ and let $A,B,C$ be the $N/3$-point DFTs of $x_{3r},x_{3r+1},x_{3r+2}$.
Show $X_k=A_k+\omega_N^kB_k+\omega_N^{2k}C_k$ for all $k$ (indices of $A,B,C$ mod
$N/3$), and deduce an $O(N\log N)$ cost when $N=3^p$.
\end{exercise}

% =====================================================================
\section{From a Spectrum to a Pitch}\label{sec:pitch}
\usedin{\code{harmonic\_product}, \code{track\_pitch} (the 10\% floor),
\code{refine\_peaks}. Slides 13 and 14.}

A hum's spectrum has a peak at every harmonic, and the tallest is often $2f_0$, an
octave too high. This lesson proves that multiplying the spectrum by compressed copies
of itself finds $f_0$, explains the guard needed for whistles, and derives the
parabola that locates a peak between bins.

\subsection{Theory}

Write $M_k=\abs{X_k}$ for the magnitude spectrum.

\begin{definition}[Harmonic product spectrum]
$P[k]=M_k\,M_{2k}\,M_{3k}$, for candidate bins $k$ with $3k$ in range.
\end{definition}

To reason about it, use an idealized hum.

\begin{definition}[Ideal harmonic spectrum]
Fix a fundamental bin $k_0\ge1$, a number of harmonics $H\ge1$, amplitudes
$a_1,\dots,a_H>0$ with $A=\max_ha_h$, and a noise level $\varepsilon\ge0$. Set
$M_{hk_0}=a_h$ for $1\le h\le H$ (the \emph{harmonic bins}) and $M_k=\varepsilon$ for
every other $k\ge1$.
\end{definition}

\begin{theorem}[The harmonic product finds the fundamental]\label{thm:hps}
In the ideal harmonic spectrum, suppose $3\le H\le5$ and $\varepsilon A^2<a_1a_2a_3$.
Then $k_0$ is the unique maximizer of $P[k]$ over $k\ge1$, even if the tallest peak is
at another harmonic.
\end{theorem}
\begin{proof}
$P[k_0]=a_1a_2a_3$. Let $k\ne k_0$. We show at most two of $k,2k,3k$ are harmonic
bins, so $P[k]\le A\cdot A\cdot\varepsilon<a_1a_2a_3$.
If $k_0\nmid k$, then $k$ itself is not a harmonic bin.
If $k=rk_0$ with $r\ge2$, then $3k=3rk_0\ge6k_0>Hk_0$, so $3k$ is not a harmonic bin.
\end{proof}

\begin{remark}[Whistles and the floor]
A whistle is nearly a pure sinusoid: $H=1$. Then $P[k_0]=a_1\varepsilon^2$, but also
$P[k_0/2]=\varepsilon\cdot a_1\cdot\varepsilon$ when $k_0$ is even, a tie that noise
breaks at random, so a subharmonic can win. The code therefore only lets a bin win if
$M_k\ge\theta\max_jM_j$ with $\theta=0.1$. If $\varepsilon<\theta A$, the admissible
bins are harmonic bins, and the subharmonics are excluded
(Exercise~\ref{ex:floor}).
\end{remark}

The winning bin is an integer, but the true pitch lies between bins. Fit a parabola
through the peak and its two neighbours.

\begin{lemma}[Vertex of a three-point parabola]\label{lem:vertex}
The parabola through $(-1,\alpha)$, $(0,\beta)$, $(1,\gamma)$ with
$\alpha-2\beta+\gamma<0$ has its maximum at
\[
\delta=\frac12\cdot\frac{\alpha-\gamma}{\alpha-2\beta+\gamma}.
\]
If moreover $\beta\ge\alpha$ and $\beta\ge\gamma$, then $\abs\delta\le\frac12$.
\end{lemma}
\begin{proof}
Write $p(u)=au^2+bu+c$. Then $c=\beta$, $a=\frac{\alpha-2\beta+\gamma}2<0$ and
$b=\frac{\gamma-\alpha}2$, and the vertex is $-b/(2a)=\delta$. For the bound:
$\abs{\alpha-\gamma}=\abs{(\beta-\gamma)-(\beta-\alpha)}\le(\beta-\alpha)+(\beta-\gamma)
=\abs{\alpha-2\beta+\gamma}$.
\end{proof}

\begin{proposition}[Exact for Gaussian peaks]
If $\log M$ near the peak is the parabola $-(u-u^\star)^2/(2\sigma^2)+\text{const}$ (a
Gaussian peak), then applying Lemma~\ref{lem:vertex} to $\log M$ at bins $k-1,k,k+1$
gives $k+\delta=u^\star$ exactly.
\end{proposition}
\begin{proof}
A parabola through three of its own points is itself, so its vertex is $u^\star$.
\end{proof}

A Hann-windowed peak is close to a Gaussian, which is why the code interpolates on
$\log$ magnitudes; zero-padding puts the three samples near the top, where the
approximation is best.

\subsection{Practice: by hand}

\begin{example}[An octave error avoided]
A low hum with a weak fundamental: $k_0=10$, $M_{10}=0.3$, $M_{20}=1$, $M_{30}=0.6$,
$M_{40}=0.3$, and $\varepsilon=0.01$ elsewhere. The tallest bin is 20, an octave too
high. But $P[10]=0.3\cdot1\cdot0.6=0.18$, $P[20]=1\cdot0.3\cdot0.01=0.003$ and
$P[5]=0.01\cdot0.3\cdot0.01=0.00003$. The condition of Theorem~\ref{thm:hps} holds:
$\varepsilon A^2=0.01<0.18$.
\end{example}

\begin{example}[Between bins]
Log-magnitudes $\alpha=-2$, $\beta=-0.5$, $\gamma=-1$ at bins $23,24,25$ give
$\delta=\frac12\cdot\frac{-1}{-2}=0.25$, so the peak is at bin $24.25$, that is
$24.25\times6.125\approx148.5$ Hz.
\end{example}

\begin{connection}
\textbf{Code connection.} \code{harmonic\_product} computes $P$ for all bins at once:
\code{hps *= spectrum[..., : h * n\_bins : h]} reads bins $0,h,2h,\dots$, which is
$M_{hk}$. \code{track\_pitch} zeroes candidates below
\code{fundamental\_floor} $=0.1$ of the loudest bin before the \code{argmax}.
\code{refine\_peaks} first climbs to the local maximum (so $\beta\ge\alpha,\gamma$ and
$\abs\delta\le\frac12$), then applies Lemma~\ref{lem:vertex} to
\code{np.log(spectrum + 1e-12)}.
\end{connection}

\subsection{Exercises}
\begin{exercise}[Hand]\label{ex:hpshand}
(a) A spectrum has $M_{12}=0.8$, $M_{24}=1$, $M_{36}=0.5$, $M_{48}=0.2$ and $0.02$
elsewhere. Which bin is tallest? Compute $P[6]$, $P[12]$ and $P[24]$; which bin does
the harmonic product choose? (b) Compute $\delta$ for
$(\alpha,\beta,\gamma)=(-1,0,-3)$ and $(-0.2,0,-0.2)$, and interpret both.
\end{exercise}

\begin{exercise}[Proof]\label{ex:floor}
In the ideal harmonic spectrum with $H\le2$, suppose $0<\varepsilon<\theta A$ and
$a_1\ge\theta A$, and admit only bins with $M_k\ge\theta A$. Show that $k_0$ is the
admissible maximizer of $P$. Why is $\varepsilon>0$ needed?
\end{exercise}

\begin{exercise}[Proof, $\star$]\label{ex:h6}
Show that Theorem~\ref{thm:hps} fails for $H=6$: find a $k\ne k_0$ for which all three
of $k,2k,3k$ are harmonic bins, and choose amplitudes that make $P[k]>P[k_0]$. What
kind of voice would do this?
\end{exercise}

% =====================================================================
\section{Notes in the Singer's Own Tuning}\label{sec:notes}
\usedin{\code{hz\_to\_midi}, \code{estimate\_detune}, \code{segment\_notes},
\code{analyze}. Slide 15.}

We now have one pitch per frame, 100 a second. Turning them into notes means taking a
logarithm and rounding, but people are rarely tuned to A440, and someone 45 cents
sharp sits right on the rounding boundary. A circular mean finds their tuning.

\subsection{Theory}

\begin{definition}[MIDI value, cents]
For a frequency $f>0$, $m(f)=69+12\log_2(f/440)$ is its (real) MIDI value; rounding
gives the nearest note. A \emph{cent} is $1/100$ of a semitone, so $m$ measures pitch
in units of 100 cents.
\end{definition}

\begin{proposition}[Logs turn ratios into differences]
$m(2^{k/12}f)=m(f)+k$, and $m(f)-m(g)=12\log_2(f/g)$. In particular transposing a
melody adds the same number to every note.
\end{proposition}
\begin{proof}
$12\log_2(2^{k/12}f/440)=k+12\log_2(f/440)$; subtract the two definitions for the
second identity.
\end{proof}

Only the fractional part of $m$ matters for tuning, and fractional parts live on a
circle: $0.99$ is next to $0.01$, not $0.98$ away. So average them as points on the
unit circle.

\begin{definition}[Detune]
For values $m_1,\dots,m_T$, let $z=\frac1T\sum_te^{2\pi im_t}$. If $z\ne0$, the
\emph{detune} is $d=\frac1{2\pi}\arg z\in(-\frac12,\frac12]$.
\end{definition}

\begin{proposition}[Invariances]\label{prop:detuneinv}
Adding integers $n_t$ to the $m_t$ does not change $d$. Adding the same real $c$ to
every $m_t$ changes $d$ to $d+c$, modulo 1.
\end{proposition}
\begin{proof}
$e^{2\pi i(m_t+n_t)}=e^{2\pi im_t}$. And $e^{2\pi i(m_t+c)}=e^{2\pi ic}e^{2\pi im_t}$,
so $z$ is multiplied by $e^{2\pi ic}$ and its angle increases by $2\pi c$.
\end{proof}

\begin{theorem}[Tuning-corrected rounding recovers the notes]\label{thm:detune}
Suppose $m_t=n_t+c+\varepsilon_t$ with integers $n_t$, a constant offset $c$, and
per-frame wobble $\abs{\varepsilon_t}\le\eta<\frac14$. Then $d\equiv c+\rho\pmod1$ with
$\abs\rho\le\eta$, and there is one integer $j$ with $\round(m_t-d)=n_t+j$ for every
$t$.
\end{theorem}
\begin{proof}
By Proposition~\ref{prop:detuneinv} we may drop the $n_t$ and rotate by $-c$: then $z$
becomes $\frac1T\sum_te^{2\pi i\varepsilon_t}$, an average of unit vectors with angles in
$[-2\pi\eta,2\pi\eta]\subset(-\frac\pi2,\frac\pi2)$. That sector is convex, so the average lies in it, and
its real part is at least $\cos(2\pi\eta)>0$, so $z\ne0$ and its angle $2\pi\rho$
satisfies $\abs\rho\le\eta$. Undoing the rotation, $d=c+\rho-j$ for some integer $j$
(bringing $d$ into $(-\frac12,\frac12]$). Then
$m_t-d=n_t+j+(\varepsilon_t-\rho)$ with $\abs{\varepsilon_t-\rho}\le2\eta<\frac12$, so it
rounds to $n_t+j$.
\end{proof}

The common integer $j$ is harmless: it is a transposition, and the matcher of
Lesson~\ref{sec:match} ignores those.

Finally, frames become notes: a run of identical rounded values shorter than 80 ms
(8 frames) is a slide between notes and is dropped, and a note equal to the previous
one only counts as new if at least 30 ms (3 frames) of silence separates them.

\subsection{Practice: by hand}

\begin{example}[A sharp singer]
Six frames hum B3 D4 A3 B3 D4 A3 (59, 62, 57, 59, 62, 57) about 45 cents sharp with a
little wobble: $m=(59.40,62.48,57.53,59.44,62.37,57.47)$. Naive rounding gives
$(59,62,\mathbf{58},59,62,57)$: the third note is wrong. The detune is
$d\approx0.448$, and $\round(m-d)=(59,62,57,59,62,57)$: all correct.
\end{example}

\begin{example}[Why not an ordinary average]
For $m=(60.95,62.05)$ the fractional parts are $0.95$ and $0.05$, whose arithmetic
mean $0.5$ says ``half a semitone off''. On the circle they are $\pm0.05$ from 0, and
the circular mean is $d=0$: in tune.
\end{example}

\begin{connection}
\textbf{Code connection.} \code{hz\_to\_midi} is $m(f)$ with NaN for unpitched
frames. \code{estimate\_detune} is literally
\begin{lstlisting}
np.angle(np.mean(np.exp(2j * np.pi * voiced))) / (2 * np.pi)
\end{lstlisting}
\code{analyze}
rounds \code{m - detune}, and \code{segment\_notes} applies the 80 ms and 30 ms rules.
The test suite hums every song 40 cents sharp to exercise Theorem~\ref{thm:detune}.
\end{connection}

\subsection{Exercises}
\begin{exercise}[Hand]\label{ex:detunehand}
(a) Find the MIDI value of 196 Hz and the distance in cents from 440 Hz to 446 Hz.
(b) Compute the detune of $m=(60.1,61.9,64.05,59.95)$ and the naive and corrected
roundings.
\end{exercise}

\begin{exercise}[Hand]\label{ex:segment}
Apply the segmentation rules to these per-frame notes (one entry per 10 ms frame, $-$
for silence): ten 62s, two 64s, nine 62s, three $-$, eight 62s. How many notes come
out, and which?
\end{exercise}

\begin{exercise}[Proof, $\star$]\label{ex:eta}
Show that the bound $\eta<\frac14$ in Theorem~\ref{thm:detune} cannot be raised to
$\frac12$: give wobbles with $\abs{\varepsilon_t}<\frac12$ for which some note is
rounded wrongly even after correction. Where does the proof break?
\end{exercise}

% =====================================================================
\section{Matching in Any Key}\label{sec:match}
\usedin{\code{rank\_songs}, \code{match\_song}, \code{DEFAULT\_TOLERANCE}. Slides 16
and 17.}

The last step compares the hummed notes with the twelve songs, ignoring key. It is
geometry: all keys of a song form a line, and the error is the distance to that line,
measured so that one wrong note cannot spoil it.

\subsection{Theory}

\begin{definition}[Key-invariant distance]
A song is a vector $s\in\Z^L$ of note numbers, and humming it $c$ semitones away gives
$s+c\one$, where $\one=(1,\dots,1)$. For a window $x\in\R^L$ of hummed notes,
\[
E(x,s)=\min_{c\in\R}\frac1L\norm{x-s-c\one}_1,
\qquad\norm v_1=\sum_i\abs{v_i}.
\]
\end{definition}

\begin{theorem}[The best shift is a median]\label{thm:median}
Let $d=x-s$. Then $g(c)=\sum_i\abs{d_i-c}$ is minimized exactly at the medians of
$d_1,\dots,d_L$ (for even $L$, every point between the two middle values).
\end{theorem}
\begin{proof}
$g$ is continuous and piecewise linear. Away from the $d_i$, its slope is
$\#\{i:d_i<c\}-\#\{i:d_i>c\}$, which increases with $c$ from $-L$ to $L$. It is
negative while fewer than half the $d_i$ lie below $c$ and positive once more than half
do, so $g$ decreases up to the median and increases after it.
\end{proof}

\begin{proposition}[The $L^2$ alternative]
$\sum_i(d_i-c)^2$ is minimized at the mean $\bar d$, and the minimum is the squared
distance from $d$ to the line $\R\one$ (orthogonal projection).
\end{proposition}
\begin{proof}
The derivative $-2\sum_i(d_i-c)$ vanishes only at $c=\bar d$, and $\bar d\one$ is the
projection $\frac{\ip d\one}{\ip\one\one}\one$.
\end{proof}

The difference matters: one wrong note drags the mean but not the median
(Example~\ref{ex:wrong}).

\begin{theorem}[A pseudometric on melodies]\label{thm:pseudo}
$E$ is symmetric, $E(x,x+c\one)=0$, and $E(x,z)\le E(x,y)+E(y,z)$ for all
$x,y,z\in\R^L$.
\end{theorem}
\begin{proof}
Symmetry: substitute $c\mapsto-c$. For the triangle inequality, take minimizers $c_1$
for $E(x,y)$ and $c_2$ for $E(y,z)$; then
$\norm{x-z-(c_1+c_2)\one}_1\le\norm{x-y-c_1\one}_1+\norm{y-z-c_2\one}_1$, and
$E(x,z)$ is at most the left side divided by $L$.
\end{proof}

\begin{corollary}[Separated songs cannot be confused]\label{cor:sep}
If two songs of length $L$ satisfy $E(s,s')>2\tau$, no window $x\in\R^L$ is within
$\tau$ of both.
\end{corollary}
\begin{proof}
Otherwise $E(s,s')\le E(s,x)+E(x,s')\le2\tau$.
\end{proof}

\begin{remark}[The songbook]
With tolerance $\tau=0.5$, Corollary~\ref{cor:sep} needs distances above 1. The nine
6-note songs are at least $1.5$ apart (Epona's Song and Requiem of Spirit are the
closest pair), so no 6-note window can match two of them. Across lengths, Requiem is
$1.17$ from a 6-note stretch of Bolero of Fire, and the corollary does not apply. In
fact exactly one hum is within tolerance of both: Bolero with its seventh note sung as
F4 instead of A4, in any key, is exactly $0.5$ from Bolero (one note off by 4, over 8
notes) and exactly $0.5$ from Requiem (its middle six notes have one note off by 3).
The code breaks ties in favour of the longer song, so it reports Bolero, and a
brute-force search over every hum within tolerance of any song finds no other
ambiguity. Exercise~\ref{ex:mixed} shows why the triangle inequality could not rule
this case out.
\end{remark}

\subsection{Practice: by hand}

\begin{example}[Zelda's Lullaby, an octave down]
$s=(71,74,69,71,74,69)$ and $x=s-12$: $d=(-12,\dots,-12)$, median $-12$, error 0. The
matcher reports ``12.0 semitones below the ocarina''.
\end{example}

\begin{example}[One wrong note]\label{ex:wrong}
Hum Zelda's Lullaby exactly but with the third note 2 semitones sharp:
$d=(0,0,2,0,0,0)$. The median is 0 and $E=2/6\approx0.33\le0.5$: a match. With the
mean $\frac13$ instead, the $L^1$ error would be
$\frac16(5\cdot\frac13+\frac53)=\frac59\approx0.56>0.5$: rejected.
\end{example}

\begin{example}[Backwards is rejected]
Reversed, Zelda's Lullaby is $(69,74,71,69,74,71)$: $d=(-2,0,2,-2,0,2)$, median 0,
$E=8/6\approx1.33$. Same notes, wrong order, far from the line.
\end{example}

\begin{connection}
\textbf{Code connection.} \code{rank\_songs} forms every window with
\code{sliding\_window\_view(played, len(melody))}, takes \code{np.median} of the
differences as the shift, and averages the absolute residuals: that is $E$, minimized
over windows. \code{match\_song} accepts the best song if $E\le0.5$. The comment on
\code{DEFAULT\_TOLERANCE} records what happens at looser settings: at 1.0, Epona's Song
backwards is accepted as Saria's Song (Exercise~\ref{ex:epona}).
\end{connection}

\subsection{Exercises}
\begin{exercise}[Hand]\label{ex:epona}
Compute $E$ between (a) Epona's Song $(74,71,69,74,71,69)$ and Requiem of Spirit
$(62,65,62,69,65,62)$; (b) Epona's Song reversed and Saria's Song
$(65,69,71,65,69,71)$. Which tolerances accept (b)?
\end{exercise}

\begin{exercise}[Proof]\label{ex:onewrong}
Suppose $x=s+c\one+\delta\,u_j$, where $u_j$ is the $j$-th standard basis vector: every
note is exact up to the key shift $c$ except note $j$, which is off by $\delta$. Show
that for $L\ge3$ the median shift is exactly $c$, and $E(x,s)=\abs\delta/L$.
\end{exercise}

\begin{exercise}[Proof, $\star$]\label{ex:mixed}
Let $B\in\Z^8$ and let $B'$ be its notes 2--7. Show that if an 8-note window $w$ has
$E(w,B)\le\tau$, then its notes 2--7, $w'$, satisfy $E(w',B')\le\frac43\tau$. Now
suppose a hum contains such a $w$ and also $E(w',R)\le\tau$ for a 6-note song $R$.
Deduce $E(R,B')\le\frac73\tau$, and compare with the Requiem--Bolero distance $1.17$ at
$\tau=0.5$.
\end{exercise}

% =====================================================================
\appendix
\section{Hints and Answers}\label{app:hints}

\paragraph{Lesson 1.} \ref{ex:missing}: $x(0)=2$, but
$x(\frac1{2f_0})=\cos2\pi+\cos3\pi=0$ and $x(\frac1{3f_0})=\cos\frac{4\pi}3+\cos2\pi=\frac12$;
$1/f_0$ is a period by Proposition~\ref{prop:harmperiod}.
\ref{ex:alias}: $44{,}100-30{,}000=14{,}100$ Hz; $50{,}000-44{,}100=5{,}900$ Hz.
\ref{ex:nyqinj}: take $n=1$; both angles $2\pi f/f_s$ and $2\pi g/f_s$ lie in
$[0,\pi]$, where cosine is one-to-one.

\paragraph{Lesson 2.} \ref{ex:sixth}: $1$, $\frac12-\frac{\sqrt3}2i$,
$-\frac12-\frac{\sqrt3}2i$, $-1$, $-\frac12+\frac{\sqrt3}2i$, $\frac12+\frac{\sqrt3}2i$;
the sum is the case $m=1$ and the sum of squares the case $m=2$ (neither divisible by
6). \ref{ex:probe4}: $1\cdot1+1\cdot(-i)+(-1)(-1)+(-1)\cdot i=2-2i$; from
$x_0=a\cos\varphi=1$ and $x_1=-a\sin\varphi=1$, $a=\sqrt2$ and $\varphi=-\pi/4$, and
$\frac{aN}2e^{i\varphi}=2\sqrt2e^{-i\pi/4}=2-2i$. \ref{ex:edgebins}:
$\cos(\pi n+\varphi)=(-1)^n\cos\varphi$, so the phase can only scale the samples.

\paragraph{Lesson 3.} \ref{ex:dft4}: $(0,2,0,2)$ and $(0,-2i,0,2i)$: equal magnitudes,
phases a quarter turn apart; the inverse of $(0,2,0,2)$ is
$\frac14(2e_1+2e_3)=(1,0,-1,0)$. \ref{ex:parseval}: $\ip{Fx}{Fy}=y^*F^*Fx$.
\ref{ex:basis}: if $\sum c_ke_k=0$, its inner product with $e_j$ is $Nc_j$; $N$
independent vectors in an $N$-dimensional space form a basis.

\paragraph{Lesson 4.} \ref{ex:layout7}: (a) bin 0 DC, bins 1--3 are $+1,+2,+3$, bins
4--6 are $-3,-2,-1$; no Nyquist bin; \code{rfft} returns $7/\!/2+1=4$ bins. (b)
$X_3=\overline{X_1}=-2-2i$ and $x=(1,2,3,4)$. \ref{ex:nyq}: $(0,0,4,0)$ for
$\varphi=0$; all zeros for $\varphi=\pi/2$, because every sample of
$\cos(\pi n+\pi/2)=-\sin(\pi n)$ is 0. \ref{ex:realseries}: in the inversion formula,
pair bins $k$ and $N-k$: by Theorem~\ref{thm:realsym},
$X_ke_k+X_{N-k}e_{N-k}=X_ke_k+\overline{X_ke_k}=2\operatorname{Re}(X_ke_k)$.

\paragraph{Lesson 5.} \ref{ex:diff}: $(Fc)_k=1-\omega^k=(0,\,1+i,\,2,\,1-i)$: DC is
erased and the Nyquist frequency doubled, a high-pass filter. \ref{ex:sunitary}: $S$
permutes the standard basis; if $Se=\lambda e$ then $e=S^Ne=\lambda^Ne$.
\ref{ex:commute}: if $AS=SA$, then column $j$ of $A$ is $Au_j=AS^ju_0=S^jAu_0$, the
first column shifted $j$ times; conversely $C=\sum_jc_jS^j$ commutes with $S$.

\paragraph{Lesson 6.} \ref{ex:hann4}: $w\odot x=(0,1,3,2)$ with DFT
$(6,-3+i,0,-3-i)$; the formula gives $5-\frac14(-4)=6$,
$(-1+i)-\frac14\cdot8=-3+i$, $-1-\frac14(-4)=0$, and $-3-i$.
\ref{ex:leak8}: $\abs{X_2}=\abs{X_3}=1/\sin\frac\pi{16}\approx5.13$ and
$\abs{X_6}=1/\sin\frac{7\pi}{16}\approx1.02$. \ref{ex:decay}: bound the numerator by
1 and apply the inequality with $y=\pi\abs{\alpha-k}/N$.

\paragraph{Lesson 7.} \ref{ex:fft4}: (a) $E=(0,2)$, $O=(0,0)$, so $X=(0,2,0,2)$.
(b) write each index in 3-bit binary and reverse the bits: $1=001\mapsto100=4$,
$3=011\mapsto110=6$, and so on. \ref{ex:radix3}: split $n=3r+j$ and use
$\omega_N^3=\omega_{N/3}$; the cost satisfies $T(N)=3T(N/3)+O(N)$.

\paragraph{Lesson 8.} \ref{ex:hpshand}: (a) the tallest bin is 24; $P[6]=0.00032$,
$P[12]=0.4$, $P[24]=0.004$; it chooses 12. (b) $\delta=-0.25$, leaning toward the
larger left neighbour; $\delta=0$, a symmetric peak. \ref{ex:floor}: the admissible
bins are $k_0$ and possibly $2k_0$; $P[k_0]=a_1a_2\varepsilon$ beats
$P[2k_0]=a_2\varepsilon^2$ exactly when $a_1>\varepsilon$; with $\varepsilon=0$ both
are 0 and nothing is decided. \ref{ex:h6}: $k=2k_0$ makes $2k_0,4k_0,6k_0$ harmonic;
$a_1=a_3=0.1$ and $a_2=a_4=a_6=1$ give $P[2k_0]=1>P[k_0]=0.01$: a voice with a weak
fundamental and strong even harmonics.

\paragraph{Lesson 9.} \ref{ex:detunehand}: (a) $69+12\log_2(196/440)\approx55.0$ (G3);
$1200\log_2(446/440)\approx23.4$ cents. (b) the offsets $+0.1,-0.1,+0.05,-0.05$ are
symmetric about 0, so $z$ is real and positive and $d=0$; both roundings give
$(60,62,64,60)$. \ref{ex:segment}: two notes, 62 and 62. The two 64-frames are shorter
than 8 frames and are dropped; the 62s on either side have no rest between them and
merge; three frames of silence then separate a second 62. \ref{ex:eta}: with $c=0$, let
one frame wobble by $+0.45$ and two by $-0.1$; then $d\approx-0.146$, and
$0.45-d\approx0.60$ rounds up while $-0.1-d\approx0.05$ rounds down. The unit vectors
no longer fit in a sector narrower than $\pi$, so their average can point outside it.

\paragraph{Lesson 10.} \ref{ex:epona}: (a) $d=(12,6,7,5,6,7)$, median $6.5$, residual
sum $9$, $E=1.5$. (b) $d=(4,2,3,4,2,3)$, median 3, $E=4/6\approx0.67$: accepted at
tolerance $1.0$, rejected at $0.5$. \ref{ex:onewrong}: $x-s=c\one+\delta u_j$ has at
least $L-1\ge2$ entries equal to $c$, more than half when $L\ge3$, so the median is
$c$. \ref{ex:mixed}: keeping six of the eight residuals can only shrink their sum,
which is at most $8\tau$; then
$E(R,B')\le E(R,w')+E(w',B')\le\tau+\frac43\tau=\frac73\tau\approx1.167$ at
$\tau=0.5$, exactly the Requiem--Bolero distance, so this argument cannot separate
them. Nothing can: the hum in the Remark on the songbook sits at exactly $0.5$ from
both songs.

\end{document}
